%% kapa: exam=1 round=25 year=2014 period=2 form=A subject=accounting subject_ko=회계학 seq=11 number=51
%% source: Q-Net 감정평가사 기출문제 (cst00309) artlSeq=5101207 file=extracted/1008/5101207/2014년 제25회 감정평가사 1차시험 기출문제/2교시(A형).hwp
%% conversion: kordoc 4.24.51 -> tools/gen1.py (auto); review=proofread-sonnet; flags=none
\begin{question}{51}
재고자산에 관한 설명으로 옳은 것은?
\begin{choices}
  \item 후속 생산단계에 투입하기 전에 보관이 필요한 경우 이외의 보관원가는 재고자산의 취득원가에 포함할 수 있다.
  \item 확정판매계약을 이행하기 위하여 보유하는 재고자산의 순실현가능가치는 계약가격에 기초하며, 확정판매계약의 이행에 필요한 수량을 초과하는 경우에는 일반 판매가격에 기초한다.
  \item 재고자산의 지역별 위치나 과세방식이 다른 경우 동일한 재고자산에 다른 단위원가 결정방법을 적용할 수 있다.
  \item 가중평균법의 경우 재고자산 원가의 평균은 기업의 상황에 따라 주기적으로 계산하거나 매입 또는 생산할 때마다 계산하여서는 아니된다.
  \item 완성될 제품이 원가 이상으로 판매될 것으로 예상하는 경우에는 해당 원재료를 순실현가능가치로 감액한다.
\end{choices}
\end{question}
\answer{2}
